\documentclass[10pt]{article}
\usepackage[a4paper,margin=24mm]{geometry}
\usepackage{amsmath,amssymb,booktabs}
\usepackage[hidelinks]{hyperref}
\newcommand{\Ai}{\operatorname{Ai}}
\newcommand{\asinh}{\operatorname{asinh}}
\newcommand{\atanh}{\operatorname{atanh}}
\newcommand{\sinc}{\operatorname{sinc}}
\title{Refined pointwise hyperspherical Bessel evaluation}
\author{}
\date{}
\begin{document}
\maketitle
\vspace{-2em}
For $K=-1,+1$ let $s=S_K(\chi)$ be $\sinh\chi$ or $\sin\chi$, respectively,
and $c=S'_K(\chi)$. Work on the folded closed interval $0\leq\chi\leq\pi/2$;
CAMB supplies the parity factor. The regular solution $u=s\Phi_\ell^\nu$ obeys
\begin{equation}
 u''=\left[\frac{\ell(\ell+1)}{s^2}-\nu^2\right]u,
 \quad u\sim A_\ell\chi^{\ell+1},\qquad
 A_\ell=\frac{\prod_{j=1}^{\ell}\sqrt{\nu^2-Kj^2}}{(2\ell+1)!!}.
 \label{eq:radial}
\end{equation}
The new routine \texttt{refined\_u} accelerates the existing pointwise Olver
fallback. It returns $u$ directly; the flat, near-flat and leading Olver paths,
coordinate folding, and source integration are unchanged.

\section{Exact representations in the degenerate limits}
\paragraph{Open geometry.}
Set $z=e^{-2\chi}$ and substitute $u_+=e^{i\nu\chi}(1-z)^{-\ell}P(z)$
into \eqref{eq:radial}. The resulting hypergeometric equation has the
terminating solution
\begin{align}
 P(z)&={}_2F_1(-\ell,-\ell-i\nu;1-i\nu;z),\qquad
 \delta=\sum_{j=1}^{\ell}\arctan(\nu/j),\\
 u&=\frac{(1-z)^{-\ell}}{\nu}\operatorname{Im}\bigl[e^{i(\nu\chi-\delta)}P(z)\bigr].
 \label{eq:jost}
\end{align}
The phase follows from the regular boundary condition (equivalently, from the
successive Darboux raising operators); the outgoing amplitude is $1$ in $u_+$.
The ratio of terms including powers of $z$ is
\begin{equation}
 \frac{T_{j+1}}{T_j}=z\frac{(\ell-j)(\ell-j+i\nu)}{(j+1)(j+1-i\nu)}.
\end{equation}
Its magnitude $q_j$ decreases with $j$. Thus, when $q_j<1$, the remaining
sum is bounded by $|T_{j+1}|/(1-q_j)$. The implementation stores
$T_j=a_j+i\nu b_j$ and evaluates \eqref{eq:jost} as
\begin{equation}
 u=(1-z)^{-\ell}\left[D\sinc(\nu D)\sum_j a_j+\cos(\nu D)\sum_j b_j\right],
 \quad D=\chi-\delta/\nu.
\end{equation}
This is continuous at $\nu=0$, where $\delta/\nu=H_\ell$ is the harmonic
number. Stable log-gamma differences and a small-$\nu$ harmonic expansion
avoid an $O(\ell)$ phase sum. Bounds on the coefficient magnitudes and phases
control the scaled imaginary tail; a separate roundoff screen rejects
ill-conditioned sums. The work limit is 40 terms. In addition to limiting
$q_0$, require $z\leq1/4$: since $q_0\geq\ell z$,
$-\ell\log(1-z)\leq q_0/(1-z)$, bounding the external amplification. The
absolute stopping tolerance is scaled by the inverse of this amplification
so that it applies to the final reduced function $u$.

\paragraph{Closed geometry.}
For $p=\nu-\ell-1=2m+\epsilon$ with $\epsilon\in\{0,1\}$, use the
Gegenbauer polynomial $s^\ell C_p^{\ell+1}(c)$, expanded about $c=0$.
In terms of $y=c^2$ and $t_0=1$,
\begin{align}
 \frac{t_{j+1}}{t_j}&=-y\frac{(m-j)(\ell+1+m+\epsilon+j)}
 {(j+1)(j+\epsilon+1/2)},\\
 u&=(-1)^m s^{\ell+1}D_0\sum_{j=0}^m t_j\times
 \begin{cases}1,&\epsilon=0,\\ c\sqrt{2(\ell+1+m)(2m+1)},&\epsilon=1,\end{cases}\\
 D_0^2&=\frac{1}{2\nu}\frac{\Gamma(\ell+1+m)}{\Gamma(\ell+3/2+m)}
 \frac{\Gamma(m+1/2)}{\Gamma(m+1)}.
\end{align}
Half-step gamma ratios are evaluated by Stirling differences. Use geometric tail and roundoff tests with an absolute floor in units of
$D_0 s^{\ell+1}$, dividing the floor by the odd-parity factor when it exceeds one.
This avoids rejecting accurate results near polynomial zeros. There is no
recurrence in multipole order. In particular $p=0$ is a single normalized $s^{\ell+1}$.

\paragraph{Compact open integral.}
For low open frequencies not covered by the short Jost series, use the
exact Mehler--Dirichlet integral (integer $\ell$)
\begin{equation}
 \Phi=\frac{N}{s^{\ell+1}}\int_0^\chi(c-\cosh t)^\ell\cos(\nu t)\,dt,
 \qquad N=\prod_{j=1}^\ell\sqrt{1+\nu^2/j^2}.
\end{equation}
With $B=\sinh^2(\chi/2)/\ell$ and $y=\sqrt\ell\sinh(t/2)/\sinh(\chi/2)$,
\begin{equation}
 u=2\sqrt B\,N\tanh^\ell(\chi/2)
 \int_0^{\sqrt\ell}\frac{(1-y^2/\ell)^\ell}{\sqrt{1+By^2}}
 \cos\bigl(2\nu\asinh(\sqrt B\,y)\bigr)\,dy .\label{eq:compact}
\end{equation}
The integrand modulus is at most $e^{-y^2}$; truncating at $y=7$ costs
at most $e^{-49}/14$ in the integral. On $\ell\geq8$, $B\leq1/4$,
$\nu^2B\leq5.25$ and $\nu\leq4.6$, use 24- and 32-point Gauss--Legendre
rules and reject if their absolute difference exceeds $2\times10^{-9}$.
This convergence check is empirical, not a certified quadrature bound.
There is no fitted tensor table. The normalization uses
$|\Gamma(1+i\nu)|^2=\pi\nu/\sinh(\pi\nu)$ and a stable Stirling difference.

\section{Explicit corrected Airy map}
Write $\lambda=\ell+1/2$, $\beta=\nu/\lambda$, $h=K/\beta^2$,
$q=s^{-2}-\beta^2$, $r=-1/(4s^2)$. The usual Liouville coordinate is
$\frac23|\zeta|^{3/2}=|\int_{\chi_t}^{\chi}\sqrt{|q|}\,d\chi|$,
with $\zeta>0$ on the evanescent side. The transformed equation is
\begin{equation}
 W_{\zeta\zeta}=(\lambda^2\zeta+\psi)W,\quad
 \psi=\frac{\zeta r}{q}+\frac{\zeta q''}{4q^2}
 -\frac{5\zeta(q')^2}{16q^3}+\frac{5}{16\zeta^2},
\end{equation}
where derivatives of $q$ are with respect to $\chi$.
Seek a corrected coordinate $g=\zeta+\lambda^{-2}B_1+\lambda^{-4}B_2$;
then $W\simeq\Ai(\lambda^{2/3}g)/\sqrt{g'}$.
The Schwarzian equation $\lambda^2g(g')^2-\{g,\zeta\}/2=\lambda^2\zeta+\psi$
gives the regular coefficient equations
\begin{align}
 B_1+2\zeta B_1'&=\psi,\\
 B_2+2\zeta B_2'&=B_1'''/2-2B_1B_1'-\zeta(B_1')^2.
\end{align}
For this potential the WKB phase primitives are elementary. Set
$v=c/\sqrt{1-\beta^2s^2}$ and $x=v^{-2}$. One way to derive them is to
expand the WKB momentum $\sqrt q F$, $F=1+\sum_{n\geq1}\lambda^{-2n}F_n$.
With $J=(v^2-1)(v^2-h)/(1-h)$, $A_0=-v(v^2-h)/(1-h)$ and
$A=A_0+J F_v/F$, the momentum equation reads
\[
 F^2=1+\lambda^{-2}[r/q+JA_v/2+A_0A/2-A^2/4].
\]
The primitives $H_n=\int(F_n/J)\,dv$ at the first two orders are
\begin{align}
 H_1(v)&=\frac{(1+h)v-5v^3/3}{8(1-h)},\quad
 H_2(v)=-\frac{v^9P(x)}{5760(1-h)^3},\\
 P(x)&=5525-9945(1+h)x+(4779+16983h+4779h^2)x^2\notag\\
 &\quad-(375+7395h+7395h^2+375h^3)x^3
 +(405h+2610h^2+405h^3)x^4 .
\end{align}
Matching to the Airy phase (or substituting into the coefficient equations)
then gives
\begin{align}
 B_1&=-H_1/\sqrt\zeta-5/(48\zeta^2),\\
 B_2&=-H_2/\sqrt\zeta-1105/(9216\zeta^5)
       -B_1^2/(4\zeta)+5B_1/(32\zeta^3).
\end{align}
These formulas use real analytic continuation across the turn. The apparent
poles cancel: write
\begin{align}
 R(x)&=\sum_{j\geq0}\frac{3(1+h+\cdots+h^j)}{2j+3}x^j,
 &\zeta&=x[(1-h)R/2]^{2/3},\\
 E(x)&=R[1-3(1+h)x/5]-1,
 &B_1\zeta^2&=5E/48,\\
 B_2\zeta^5&=(P R^3-5525+750E-125E^2)/46080 .
\end{align}
The numerators vanish to order 2 and 5 respectively. In the local disk
$|x|\max(1,|h|)\leq0.2$, evaluate divided series in the scaled $x$;
coefficients of the second numerator are generated in a symmetric polynomial
basis by \texttt{tests/generate\_refined\_airy.py}. Outside the disk the
closed forms and their analytical derivatives are well conditioned.

With $a_0=\lim_{\chi\to0}[\int_\chi^{\chi_t}\sqrt q\,dt+\log\chi]$,
match the regular origin amplitude using
\begin{equation}
 u\simeq M(\lambda^{2/3}\zeta/q)^{1/4}
 \frac{\Ai(\lambda^{2/3}g)}{\sqrt{g'}},\quad
 M=2\sqrt\pi A_\ell\exp\!\left(\lambda a_0-\frac{H_1(1)}{\lambda}
 -\frac{H_2(1)}{\lambda^3}\right).
\end{equation}
Omit $H_2$ when using only $B_1$; $a_0=\log(2/\sqrt{\beta^2-K})-\beta T_K^{-1}(1/\beta)$,
where $T_{-1}^{-1}=\arctan$, $T_1^{-1}=\atanh$.
The code cancels the large gamma and action terms analytically and retains
small Stirling remainders. It works in the scaled coordinate
$\tau=\lambda^{2/3}\zeta$ and preserves the corrected action in the Airy tails.

\section{Selection, error control and validation}
A useful common expansion parameter is
$\mu=\lambda(1-h)/(1+|h|)^{3/2}$: it approaches $\lambda$ near flat space,
$\nu$ for small open $\nu$, and $(\nu-\lambda)/\sqrt2$ near the closed
sectoral limit. Use the first corrected map for $\mu\geq16$, both
corrections for $4\leq\mu<16$, and exact series or \eqref{eq:compact} otherwise.
Short exact series take precedence when economical. In the open region
$\mu<4$, $\ell\geq8$ implies $\nu^2<16(8.5)^2/((8.5)^2-16)$, so
$B\leq1/4$ fits the compact box; for $B\geq1/4$, $q_0\leq1/(16B^2)\leq1$.
This is a coverage argument, not a proof of the numerical error of the maps.
Term bounds, roundoff screens, positive map Jacobian and finite-value checks
reject unsafe cases. The original CAMB fallback remains available.

The exact series have analytical truncation bounds (the floating-point
roundoff screens are heuristic). The Airy order thresholds and quadrature
convergence criterion are validated choices, not rigorous continuous-domain
forward-error certificates. Near zeros the relevant accuracy is absolute or
peak-normalized, not relative. No new hyperspherical recurrence is used.
On a retained 64,323-point positive-frequency stress grid (both curvature
signs and flat space), the complete dispatcher had no errors above
$2\times10^{-4}$; the new paths directly accepted 20,678 of 20,879 curved
fallback calls with $\ell>2$. An independent turning-point grid and a
compact-domain boundary grid, the latter extending to $\ell\simeq10^6$,
also passed the corresponding peak/local-scale tests. In a standalone
Fortran harness with the actual CAMB numerical routines, the median
CPU time per point was 2.35 versus 0.52 microseconds (original versus
candidate, identical GNU Fortran optimization flags). These are synthetic
pointwise timings, not cosmological speedups. The absolute error floors also
avoid 266 rejected polynomial evaluations in 3,024 targeted closed-zero tests.

A full single-thread CAMB comparison against the matching baseline at
$\ell_{\max}=1200$ and five values of $\Omega_K$ between $-0.05$ and $0.05$
found a largest
per-spectrum peak-normalized difference of $4.61\times10^{-5}$, and largest
source-normalized scalar transfer difference of $3.98\times10^{-5}$.
Flat results were identical; whole-model runtimes were essentially unchanged.
A mode-wide closed gate was also tested: although it suppressed the old radial
switch discontinuity, it changed the strongly closed transfer data by up to
$2.66\times10^{-3}$ of the source maximum, so it is not included here.
These comparisons establish agreement with the baseline, not absolute
cosmological accuracy. Retain the existing leading map and safety paths
pending higher-accuracy cosmological reference tests.
\end{document}
